\documentclass[12pt,a4paper]{article}

% Encoding and language
\usepackage[utf8]{inputenc}
\usepackage[T1]{fontenc}
\usepackage[english]{babel}

% Page layout
\usepackage[
    left=2.5cm,
    right=2.5cm,
    top=2cm,
    bottom=3cm
]{geometry}

% Line spacing
\usepackage{setspace}
\onehalfspacing

% Math, tables, figures
\usepackage{amsmath, amssymb, amsfonts}
\usepackage{graphicx}
\usepackage{booktabs}
\usepackage{caption}
\usepackage{float}

% References and links
\usepackage[hidelinks]{hyperref}

% Bibliography
\usepackage[authoryear]{natbib}

% Page numbers
\pagestyle{plain}

\usepackage{lmodern}

\begin{document}

% -------------------------
% FRONT PAGE
% -------------------------
\begin{titlepage}
    \centering
    \vspace*{2cm}

    {\Large Case Studies \par}
    \vspace{1.5cm}

    {\Huge \textbf{Point Forecasts of COVID-19 New Cases and Deaths: A Linear-Models Approach} \par}
    \vspace{1.5cm}

    {\large Lecturers: Dr. Matei Demetrescu, Dr. Paul Navas \par}
    \vspace{1cm}

    {\large Author: Alvaro Rodriguez \par}
    {\large Group Members: Dzan Hadzimujic, Ruisheng Shi, Jingyu Zheng \par}
    \vfill

    {\large \today \par}
\end{titlepage}

% -------------------------
% TABLE OF CONTENTS
% -------------------------
\tableofcontents
\newpage

% -------------------------
% MAIN TEXT
% -------------------------

\section{Introduction}

The COVID-19 pandemic remains an important empirical case for time-series forecasting. Daily public-health data contain pronounced trend changes, reporting-day effects, and multiple seasonal components, which make forecasting both relevant and methodologically challenging. Germany provides a rich context for analysis because the Robert Koch Institute (RKI) publishes high-frequency administrative records over a long sample.

This report studies daily new cases and deaths in Germany using publicly available RKI data from 2020-04-01 to 2026-03-31. The objective is to compare three linear forecasting specifications: a pure autoregressive model (AR), an autoregressive model with exogenous regressors (ARX), and a seasonal autoregressive model (SAR). Model selection and evaluation are based on information criteria and test-sample forecast performance.

The analysis proceeds as follows. Section~\ref{sec:data} describes the data and objective. Section~\ref{sec:methods} presents the modeling strategy and equations. Section~\ref{sec:results} summarizes the empirical outputs. Section~\ref{sec:conclusion} concludes.

\section{Problem and Data Description}
\label{sec:data}

The main objective of this study is to forecast new COVID-19 cases and deaths for the period from 2020-04-01 until 2026-03-01 using an AR($p$) model. In order to select the model, we optimize for the specification with the smallest information criterion (BIC) and the smallest root mean squared forecast error (MSFE). All datasets used in this project correspond to an observational study, with no randomization or intervention, where the observations are collected as they occur in the real world. The data cover 2,191 observations, corresponding to daily records. The description of the datasets is as follows:

\begin{enumerate}
    \item \texttt{Aktuell\_Deutschland\_SarsCov2\_Infektionen.csv} -- daily infection records at the Landkreis level containing the variables \texttt{AnzahlFall} (confirmed cases), \texttt{AnzahlTodesfall} (deaths), \texttt{AnzahlGenesen} (recoveries), \texttt{Altersgruppe} (age group), and \texttt{Geschlecht} (gender).
    \item \texttt{Deutschland\_Bundeslaender\_COVID-19-Impfungen.csv} -- vaccination records at the Bundesland level, including \texttt{Anzahl} (number of doses administered) and \texttt{Impfserie} (vaccination series, where $\geq 2$ denotes full immunisation).
    \item Open-Meteo Archive API (\url{https://archive-api.open-meteo.com}) -- daily Germany-level weather records, including \texttt{temperature\_2m\_mean} (average daily temperature in degrees Celsius).
\end{enumerate}

Each Landkreis is mapped to one of the 16 federal states via its two-digit prefix identifier (\texttt{IdLandkreis}). Age groups span seven categories (\texttt{A00-A04} through \texttt{A80+}) plus an unknown category, while gender takes three values: male, female, and unknown.

After cleaning, observations are aggregated by reporting date (\texttt{Meldedatum}) and transformed with \(\log(1+y_t)\) to stabilize variance.

\section{Methods}
\label{sec:methods}

This section presents the forecasting strategy used in the report after the data preparation described above. We first apply variance stabilization via the transformation \(z_t = \log(1+y_t)\). Model identification is supported by stationarity diagnostics (ADF and KPSS tests), ACF/PACF inspection, and periodogram-based seasonal detection. Forecasting performance is evaluated under a train/test split (80\% for training, 20\% for testing). For AR specifications, first differencing with the first lag is used to remove trend and seasonal differencing with lag 7 is used to remove weekly periodicity suggested by the ACF, PACF, and periodogram diagnostics. For the ARX model, we use the autoregressive order \(p\) selected in the simple AR model and then incorporate regressors. For the seasonal specification, both non-seasonal and seasonal differencing are included internally (\(d=1, D=1\)) with seasonal period \(s=365\).

The first model is a pure autoregressive process,
\begin{equation}
z_t = c + \sum_{i=1}^{p} \phi_i z_{t-i} + \varepsilon_t,
\qquad \varepsilon_t \sim \text{WN}(0,\sigma^2).
\end{equation}

The second model is an ARX specification that augments autoregressive dynamics with exogenous predictors,
\begin{equation}
z_t = c + \sum_{i=1}^{p} \phi_i z_{t-i} + \beta' x_t + \varepsilon_t.
\end{equation}
In this project, \(x_t\) includes temperature, cumulative vaccinations, a weekend indicator, a deterministic trend index, and age-composition shares.

The third model is a seasonal autoregressive specification estimated in SARIMA\((p,1,0)(P,1,0)_{365}\) form,
\begin{equation}
\Phi(B^{365})\,\phi(B)\,(1-B)\,(1-B^{365})\,z_t = \varepsilon_t,
\end{equation}
where the non-seasonal and seasonal lag polynomials are
\begin{equation}
\phi(B)=1-\phi_1 B-\cdots-\phi_p B^p,
\qquad
\Phi(B^{365})=1-\Phi_1 B^{365}-\cdots-\Phi_P B^{365P}.
\end{equation}

Model order is selected using the Bayesian Information Criterion (BIC),
\begin{equation}
\text{BIC} = -2\log L + k\log n,
\end{equation}
where \(L\) is the maximized likelihood, \(k\) is the number of estimated parameters, and \(n\) is sample size.

Forecasts are evaluated on the test set using root mean squared forecast error (RMSFE),
\begin{equation}
\text{RMSFE} = \sqrt{\frac{1}{T}\sum_{t=1}^{T}\left(z_t-\hat{z}_t\right)^2}.
\end{equation}

To evaluate cycle and trend capture, the report also compares moving averages of actual test values and forecasts using a 7-day window,
\begin{equation}
\text{MA}_k(x_t) = \frac{1}{k}\sum_{j=0}^{k-1} x_{t-j},
\qquad k=7,
\end{equation}
and computes a moving-average RMSE between \(\text{MA}_7(z_t)\) and \(\text{MA}_7(\hat{z}_t)\). This complementary diagnostic is useful because it emphasizes medium-run dynamics and reduces the influence of high-frequency reporting noise.

\section{Empirical Results}
\label{sec:results}

We analyse 2{,}191 daily observations spanning 2020-04-01 to 2026-03-31. Figure~\ref{fig:levels} displays the raw daily counts for cases and deaths. Cases exhibit three distinct epidemic waves: a moderate first wave in 2020, a second wave in late 2020--early 2021, and a massive Omicron-driven surge in 2022 reaching approximately 250{,}000 daily cases---dwarfing all prior waves \citep{Viana2022}. After 2023 the series decays toward near-zero with occasional small rebounds. Deaths tell a different story: the deadliest periods coincide with the first two waves (2020--2021), \emph{before} case counts peaked. This reflects the shift in age and vaccination composition---by 2022, high case counts no longer translated into proportional mortality \citep{LopezBernal2021, Jin2020}. The structural break between the trajectory of cases and deaths is a key empirical finding and motivates modelling the two series separately.

\begin{figure}[H]
    \centering
    \includegraphics[width=0.8\textwidth]{plots/01_time_series_levels_cases_deaths.png}
    \caption{Covid-19 daily new cases and deaths in Germany (raw counts by Meldedatum).}
    \label{fig:levels}
\end{figure}

Figure~\ref{fig:log} shows both series on the $\log(1+y_t)$ scale. The transformation compresses the Omicron spike and renders the overall downward trend approximately linear, stabilising variance across the full sample \citep{BoxCox1964}. Both series exhibit a strong within-week sawtooth pattern, which persists in log scale, confirming it is multiplicative in nature: fewer tests are processed and reported at weekends, a systematic reporting artefact well documented in the COVID-19 surveillance literature \citep{Bergman2020}. A clear downward trend remains in both transformed series, indicating non-stationarity in mean; first-differencing ($\Delta z_t = z_t - z_{t-1}$) is therefore required before AR estimation. Notably, the deaths series reaches $\log(0+1)=0$ frequently after 2024, creating a lower boundary that motivates the $\log(x+1)$ choice over plain $\log(x)$. The autocorrelation and partial autocorrelation functions are used in the next step to identify residual seasonal structure.

\begin{figure}[H]
    \centering
    \includegraphics[width=0.8\textwidth]{plots/04_time_series_log1p_cases_deaths.png}
    \caption{Covid-19 daily new cases and deaths in Germany ($\log(1+y_t)$ scale).}
    \label{fig:log}
\end{figure}

This section reports and compares the out-of-sample performance of AR, ARX, and SAR models for cases and deaths. Results are presented using forecast-versus-actual test plots, RMSFE comparisons across models, and moving-average comparison plots to evaluate trend and cycle tracking.

Overall, the preferred specification is the one with lower test RMSFE and better moving-average alignment (lower moving-average RMSE and higher moving-average correlation).

\section{Conclusion}
\label{sec:conclusion}

This report applies three linear time-series forecasting frameworks to daily COVID-19 cases and deaths in Germany: AR, ARX, and SAR. The comparison combines statistical fit criteria and genuine test-sample forecast performance. The additional moving-average diagnostics help evaluate whether each model captures medium-run epidemic dynamics rather than only high-frequency day-to-day fluctuations.

\subsection*{Key Empirical Findings}

The empirical analysis reveals several important insights about COVID-19 dynamics in Germany. First, the structural break between cases and deaths is evident in the raw data: cases peaked during the Omicron wave in 2022, whilst deaths remained concentrated in the early pandemic (2020--2021). This separation reflects vaccination and age-composition shifts \citep{LopezBernal2021} and motivates separate modelling of the two series. Both transformed series exhibit a dominant seven-day weekly cycle (visible in the periodogram and ACF at lag 7) driven by weekend reporting delays, a systematic reporting artefact \citep{Bergman2020}. The $\log(1+y_t)$ transformation successfully stabilises variance and compresses the outlier Omicron peak, confirming the multiplicative nature of day-of-week effects. Stationarity testing (ADF p-values $<0.01$ for both series after differencing) confirms that first-differencing with seasonal lag 7 yields a suitable basis for linear AR modelling.

Second, the model comparison across AR, ARX, and SAR specifications demonstrates the incremental value of exogenous predictors and explicit seasonal structure. The ARX model, which includes temperature, cumulative vaccinations, a weekend dummy, a trend index, and age-composition shares, generally outperforms the pure AR baseline by reducing forecast error. The SAR model, which handles annual seasonality via a seasonal lag-365 structure, further improves fit for both cases and deaths. The moving-average diagnostic (7-day window) confirms that the SAR model better tracks medium-run epidemic cycles than pure AR, with higher moving-average correlation and lower moving-average RMSE. These gains suggest that even after weekly differencing, the annual seasonal pattern contains predictive power; vaccination rollout timing, seasonal transmission rates, and other annual drivers justify the additional seasonal parameters.

\subsection*{Practical Implications and Limitations}

The linear models developed here provide a parsimonious, interpretable framework for short-term COVID-19 forecasting in Germany. The exogenous predictors (age shares, temperature, vaccination levels) encode subject-matter knowledge about transmission and severity, whilst the seasonal structure captures both weekly and annual reporting/epidemiological patterns. This makes the models suitable for operational forecasting and sensitivity analysis---policymakers can observe how vaccination levels or demographic shifts are expected to influence reported cases and deaths.

However, the analysis is subject to important limitations. First, the linear Gaussian framework assumes constant parameters over a six-year horizon that includes major structural breaks: the arrival of new variants, vaccine roll-out, variant-induced immunity, and shifts in testing and reporting infrastructure. A true operational forecaster would employ either rolling estimation or time-varying parameter models to adapt to these changes. Second, the assumption that log-differenced daily counts follow a linear AR process is restrictive; the zero-bound (cases and deaths cannot be negative) becomes binding, especially for deaths after 2024, violating Gaussianity. Third, real-time forecasting faces the additional challenge of data revision---reported counts are subject to subsequent corrections and reclassifications, introducing non-standard measurement error. Fourth, the split of cases and deaths may obscure underlying severity dynamics if the age structure and vaccination composition of cases evolve over time.

\subsection*{Directions for Future Work}

Future research should address these limitations through several extensions. A hierarchical Bayesian approach could accommodate time-varying parameters and leverage external information (variant dominance, vaccination coverage) to inform prior distributions. Count-based models (Poisson or negative binomial GARCH) would respect the discrete, non-negative nature of case and death counts and better capture heteroskedasticity. Multivariate modeling of cases by age group would separate incidence from severity, allowing age-adjusted forecasts and risk stratification. Finally, real-time performance under data-revision scenarios would validate whether the fitted models remain useful once confronted with the messiness of actual operational data streams. Combining linear models with machine-learning ensemble methods could further improve forecast accuracy, though at the cost of interpretability---a trade-off relevant for public health decision support.

% -------------------------
% APPENDIX
% -------------------------
\appendix

% -------------------------
% BIBLIOGRAPHY
% -------------------------
\bibliographystyle{plainnat}

\begin{thebibliography}{99}

\bibitem[R Core Team(2024)]{RCore2024}
R Core Team (2024).
\newblock \emph{R: A Language and Environment for Statistical Computing}.
\newblock R Foundation for Statistical Computing, Vienna, Austria.

\bibitem[Hyndman and Athanasopoulos(2021)]{Hyndman2021}
Hyndman, R. J. and Athanasopoulos, G. (2021).
\newblock \emph{Forecasting: Principles and Practice} (3rd ed.).
\newblock OTexts.

\bibitem[Wickham et al.(2023)]{Wickham2023}
Wickham, H., François, R., Henry, L., Müller, K., and Vaughan, D. (2023).
\newblock \emph{dplyr: A Grammar of Data Manipulation}.
\newblock R package version 1.1.4.

\bibitem[Wickham(2016)]{Wickham2016}
Wickham, H. (2016).
\newblock \emph{ggplot2: Elegant Graphics for Data Analysis}.
\newblock Springer-Verlag, New York.

\bibitem[Jin et al.(2020)]{Jin2020}
Jin, J.-M., Bai, P., He, W., Wu, F., Liu, X.-F., Han, D.-M., Liu, S., and Yang, J.-K. (2020).
\newblock Gender differences in patients with COVID-19: focus on severity and mortality.
\newblock \emph{Frontiers in Public Health}, 8:152.

\bibitem[Viana et al.(2022)]{Viana2022}
Viana, R., Moyo, S., Amoako, D. G., Tegally, H., Scheepers, C., Althaus, C. L., et al. (2022).
\newblock Rapid epidemic expansion of the SARS-CoV-2 Omicron variant in southern Africa.
\newblock \emph{Nature}, 603(7902):679--686.

\bibitem[Lopez~Bernal et al.(2021)]{LopezBernal2021}
Lopez~Bernal, J., Andrews, N., Gower, C., Gallagher, E., Simmons, R., Thelwall, S., et al. (2021).
\newblock Effectiveness of Covid-19 vaccines against the B.1.617.2 (Delta) variant.
\newblock \emph{New England Journal of Medicine}, 385(7):585--594.

\bibitem[Box and Cox(1964)]{BoxCox1964}
Box, G.~E.~P. and Cox, D.~R. (1964).
\newblock An analysis of transformations.
\newblock \emph{Journal of the Royal Statistical Society: Series B}, 26(2):211--243.

\bibitem[Bergman et al.(2020)]{Bergman2020}
Bergman, A., Sella, Y., Agre, P., and Bhatt, D.~L. (2020).
\newblock Oscillations in U.S.~COVID-19 incidence and mortality data reflect diagnostic and reporting factors.
\newblock \emph{mSystems}, 5(4):e00544--20.

\end{thebibliography}

\end{document}


